\section{Fluid Dynamics and Equilibrium}\label{sec:fluid}
Unlike traditional queueing systems where job sizes are fixed, LLM inference faces a unique challenge: KV cache memory grows dynamically as prompts progress through decode stages. Each additional generated token increases the prompt's memory footprint, making it difficult to predict when total memory usage will exceed capacity $C$ and trigger costly evictions. These evictions waste the partial computation already invested in interrupted prompts, degrading both throughput and latency.

We develop a fluid model to characterize the system's equilibrium state---a balanced configuration where arrival rates match completion rates and memory usage stabilizes. This equilibrium serves two purposes. First, it quantifies the maximum sustainable throughput $\throughput^*$ under memory constraint $C$, establishing a benchmark for any scheduling policy. Second, it reveals the load-balancing principle that guides our algorithm design (Section~\ref{sec:wait}): maintain the system close to equilibrium to prevent evictions. When prompts distribute evenly across processing stages, memory consumption remains predictable and overflow risk minimizes.

We first analyze a system with a single prompt type to build intuition about equilibrium dynamics, then extend to multiple heterogeneous types. The equilibrium provides a theoretical benchmark; our WAIT algorithm (Section~\ref{sec:wait}) achieves asymptotic optimality by approaching this benchmark through threshold-based admission control that keeps the system near the balanced state.

\subsection{Single-Type Fluid Model}
We begin with a system handling prompts of a single type \(j\). Each prompt's memory usage depends on its processing stage: \(l_j\) tokens in the prefill stage ($s=0$) and \(l_j + s\) tokens in the \(s\)-th decode stage (\(s = 1, \dots, l_j'\)). When the prompt completes stage $s = l_j'$, it finishes decoding and its KV cache is cleared from GPU memory.

In equilibrium, we assume the system maintains \(n_j^*\) active prompts distributed evenly across all $(l_j'+1)$ stages, with \(n_j^*/(l_j'+1)\) prompts at each stage. This balanced distribution is crucial: if prompts accumulate at later decode stages---where memory consumption is highest ($l_j + l_j'$ tokens)---total usage may exceed capacity $C$, forcing evictions that waste partial computation. By spreading prompts evenly, we minimize overflow risk while maximizing throughput. This equilibrium model allows us to derive the system's fundamental capacity limit.

Before deriving the equilibrium parameters, we identify when the system becomes fundamentally unstable---regardless of the scheduling policy used. Intuitively, each prompt requires total processing time $d_1 (l_j' + 1)(l_j + l_j'/2)$ to load its KV cache across all $(l_j'+1)$ stages. When the arrival rate exceeds the inverse of this processing time, the queue grows unbounded.

\begin{proposition}\label{prop:large_rate}
When the condition $\lambda_j d_1 (l_j' + 1) \left( l_j + \frac{l_j'}{2} \right) \geq 1$ is satisfied, the system becomes unstable in the sense that the expected average latency under any scheduling policy with deterministic arrival rate \( \lambda_j \) grows linearly with time. Formally,
$$
\mathbb{E}[\textbf{Latency}^{(T,\pi)}] = \Omega(T) \quad \text{as} \quad T \to \infty.
$$
\end{proposition}

This condition represents the fundamental capacity limit: arrivals outpace completions, causing latency to grow linearly with time. The proof constructs an equilibrium batch and shows that even in this balanced state, arrivals exceed completions when the condition holds (Appendix~\ref{appendix:proof_large_rate}). We therefore assume $\lambda_j d_1 (l_j' + 1) \left( l_j + \frac{l_j'}{2} \right) <1$ for the remainder of this section, ensuring the system can reach equilibrium.

% \arc{anyway, prove that when $n_j^*$ has not positive solution, then latency cannot be controlled. ,there is an instance such that $\ex{}{\text{Latency}^T}\to\infty$ as $T\to\infty.$}

The total memory usage in equilibrium depends on the distribution of prompts across stages:
\begin{equation}
    M_j^* = \frac{n_j^*}{l_j^{\prime}+1}\sum_{s=0}^{l_j^{\prime}}(l_j+s) = n_j^*\prn{l_j+\frac{l_j'}{2}},
\end{equation}
where we sum over all $(l_j'+1)$ stages, each containing $n_j^*/(l_j'+1)$ prompts. The factor $(l_j + l_j'/2)$ represents the average memory per prompt across its lifecycle, from prefill ($l_j$) to final decode stage ($l_j + l_j'$).

To find $n_j^*$, we use the equilibrium condition: arrivals must equal completions. Each iteration processes one stage for all $n_j^*$ prompts, taking time $(d_0 + d_1 M_j^*)$. During this time, $\lambda_j(d_0 + d_1 M_j^*)$ new prompts arrive. For equilibrium, this must equal the $n_j^*/(l_j'+1)$ prompts that complete:
$$ \underset{\text{Time per iteration}}{\underbrace{(d_0+d_1M_j^{*} )}}\lambda_j =\left(d_0+d_1 n_j^*\left( l_j +\frac{l_j^{\prime}}{2} \right) \right)\lambda_j= \frac{n_j^*}{l_j^{\prime}+1}. $$

Solving for $n_j^*$ yields:
$$n_j^* = \frac{d_0\lambda_j(l_j^{\prime}+1)}{1-d_1\lambda_j(l_j^{\prime}+1)(l_j+\frac{l_j^{\prime}}{2})}.$$

Substituting this into the memory equation, the equilibrium memory equals:
\begin{equation}
    \label{eq:memory_single}
M_j^* = n_j^*\left( l_j+\frac{l_j^{\prime}}{2} \right)
=  \frac{d_0\lambda_j(l_j^{\prime}+1)\left( l_j+\frac{l_j^{\prime}}{2} \right) }{1-d_1\lambda_j(l_j^{\prime}+1)(l_j+\frac{l_j^{\prime}}{2})}.
\end{equation}

This equilibrium batch size balances arrivals and completions. The denominator $\left(1 - d_1\lambda_j(l_j'+1)(l_j + l_j'/2)\right)$ represents the fraction of iteration time available for processing---as this approaches zero, batch size grows unbounded (approaching instability).
% \begin{comment}
% When the available memory is bigger than the equilibrium memory $M_j^*$, we can form a batch with a bigger  batch size. For example, each stage has $\frac{n_j^*}{l_j^\prime+1}+1$ prompts, and the memory $M_j^{+}$ is that
% $$M_j^{+} = \left(\frac{n_j^*}{l_j^\prime+1}+1\right)\sum_{k=0}^{l_j^{\prime}}(l_j+k)=(n_j^*+l_j^\prime+1)\left( l_j + \frac{l_j^\prime}{2} \right) = M_j^* + (l_j^\prime+1)\left( l_j + \frac{l_j^\prime}{2} \right)  $$
% So the arrival and completion ratio is that 
% $$ 
% \frac{\text{Arrival}}{\text{Completion}} 
% = \frac{\left(d_0+d_1M_j^*+d_1 (l_j^\prime+1)\left( l_j + \frac{l_j^\prime}{2} \right)\right)\lambda_j}{\frac{n_j^*}{l_j^\prime+1}+1} 
% = \frac{n_j^*+d_1 \lambda_j(l_j^\prime+1)^2\left( l_j + \frac{l_j^\prime}{2} \right)}{n_j^*+l_j^\prime+1} < 1 $$
% Where the last inequality comes from that $d_1 \lambda_j(l_j^\prime+1)\left( l_j + \frac{l_j^\prime}{2} \right)<1$. Similarly, when the available memory is smaller than equilibrium memory $M_j^*$, we have to form a batch with a smaller batch size. For example, each stage has $\frac{n_j^*+1}{l_j^\prime} -1$ prompts, and the memory $M_j^{-}$ is that
% $$
% M_j^{-}
% = \left( \frac{n_j^*}{l_j^\prime+1}-1 \right) \sum_{k=0}^{l_j^\prime}(l_j+k) = M_j^* - (l_j^\prime+1)\left(l_j + \frac{l_j^\prime}{2} \right)
% $$
% So the arrival and completion ratio is that 
% $$
% \frac{\text{Arrival}}{\text{Completion}} 
% = \frac{n_j^* - d_1\lambda_j(l_j^\prime+1)^2\left( l_j + \frac{l_j^\prime}{2} \right) }{ n_j^* - (l_j^\prime+1)  }>1
% $$
% \end{comment}
The average throughput at equilibrium is the number of decode tokens generated per iteration divided by the time consumption for one iteration. In each iteration, $n_j^*/(l_j'+1)$ prompts complete (those at stage $l_j'$), each generating $l_j'$ decode tokens:
\begin{equation}\label{eq:throughput_fluid}
\throughput_j^* = \frac{n_j^* \cdot l_j'}{(l_j'+1)(d_0+d_1n_j^*\prn{l_j+\frac{l_j'}{2}})} = \lambda_j l_j',
\end{equation}
where the equality follows from the equilibrium condition in line 46. This shows how arrival rate, memory, and throughput relate in equilibrium, where arrivals equal completions.

% \gan{I should add more explanation. Why we introduce this class.}
% Moreover, we define a class of scheduling policies \(\Pi\) that includes non-predictive policies without \textit{preemption}. Here, \textit{non-predictive} means the policy uses only the realized sample path up to the current time, without knowing future information. The term \textit{preemption} describes the action where the policy clears the memory used by an ongoing prompt request and re-queues it to the front of the waiting prompts, usually because of limited memory. Policies in \(\Pi\) avoid preemption because it causes problems. For instance, preemption interrupts the ongoing work, requiring the system to restart or reload data for that prompt, which takes extra time and uses more computing resources. This delay in finishing the preempted prompt also lowers the system’s overall efficiency by repeating tasks unnecessarily. As an example, a non-predictive policy that reserves memory of \(l_{\max} + l_{\max}'\) for each prompt—where \(l_{\max} + l_{\max}'\) is the maximum memory one prompt can need—avoids preemption by always having enough memory, making it a policy in \(\Pi\) (though in this way, it may utilize the memory inefficiently).
% \gan{cite some example of scheduling policies without Preemption.}

% \begin{proposition}\label{prop::online policy expected throughput upper bound}
% Suppose the memory capacity is not smaller than the equilibrium memory capacity $M^*_j$, consider the expected upper bound of any online policy $c\in \Pi$ in the stochastic scenario, compared with the throughput of the steady system, we have that
% $$\ex{}{\throughput^{c}(T)}\leq \throughput^*= \lambda_j(l_j^\prime+1)$$
% \end{proposition}
% So for any policy without \textbf{Preemption}, the expectation of average throughput is bounded by $\throughput^*$. Detailed proofs are provided in Appendix~\ref{appendix::proof of online policy expected throughput upper bound}
% \begin{comment}The time cost for processing $K\cdot n_j^*$ prompts of type $j$ of static system is $Kd_0+d_{c}KM^*_j$. Under any online policy $c\in \Pi$, the time cost is $\tilde{N}\cdot d_0+d_{c}KM^*_j$ under stochastic scenario, where $\tilde{N}$ is the number of batch iterations. The second term is the same since this term only corresponds to the total memory of all the processed prompts, which is $KM_j^*$. For the first term, it only corresponds to the number of batch iterations to complete the inference of all the prompts. Since the accumulated total memory usage is $KM_j^*$, and the biggest batch size is $M_j^*$, we can easily derive that $\tilde{N}\geq K$. So we have that
% $$\ex{}{\text{Average Throughput}^{c}(T)}\leq \text{Average Throughput}^*= \lambda_j(l_j^\prime+1) $$
% \end{comment}




\subsection{Multiple-Type Fluid Model}
The single-type analysis provides intuition, but real systems handle multiple prompt types simultaneously, each with different input lengths $l_j$ and output lengths $l_j'$. This heterogeneity introduces a key challenge: types compete for the same memory capacity $C$. Admitting a long prompt (large $l_j'$) may force evicting multiple short prompts, wasting their partial computation---a cascading failure we analyze in Section~\ref{sec:wait}. We now extend the equilibrium model to capture this multi-type interaction.

We assume the system maintains $n_j^*$ active prompts of each type $j \in \{1, \dots, m\}$ in equilibrium. Following the derivation of \eqref{eq:memory_single}, the total memory occupancy equals:
\begin{equation}
\label{eq:memory_multi1}
    M^* = \sum_{j=1}^m n_j^*\left(l_j + \frac{l_j^\prime}{2}\right)
\end{equation}
with:
$$ \underset{\text{Time per iteration}}{\underbrace{(d_0+d_1M^*)}}\lambda_j(l_j'+1) = n_j^*, \;\forall j \in \{1,2,\cdots,m\}$$
since the arrivals during each iteration should equal the number of completions in equilibrium, and there are $(l_j'+1)$ stages total. Substituting the above equations into \eqref{eq:memory_multi1} yields
\begin{equation}
\label{eq:memory_multi2}
M^* = \frac{d_0\sum_{j=1}^m \lambda_j (l_j^{\prime}+1)\left(l_j + \frac{l_j^{\prime}}{2} \right)}{1-d_1\sum_{j=1}^m \lambda_j (l_j^{\prime}+1)\left(l_j + \frac{l_j^{\prime}}{2} \right) } 
\end{equation}
Aggregating the results above, we get the processing time per iteration in equilibrium:
\begin{equation}    \label{eq:time_fluid}
    \begin{aligned}
            &\text{Processing Time}^* \\&= d_0 + d_1 M^*  \\
                  &= \frac{d_0}{1 - d_1 \sum_{j=1}^m \lambda_j (l_j' + 1)(l_j + \tfrac{1} {2}l_j')},
    \end{aligned}
\end{equation}
as well as the average throughput:
\begin{equation}
    \label{throughput_fluid_multiple}
    \begin{aligned}
        &\throughput^*
        = \sum_{j=1}^m \frac{n_j^* \cdot l_j^{\prime}}{(l_j^{\prime}+1) \cdot (d_0 + d_1 M^*)}
        = \sum_{j=1}^m \lambda_j l_j^{\prime}
    \end{aligned}
\end{equation}
% \begin{equation}
%             \throughput^* 
%         = \frac{1}{d_0}\left(\sum_{j=1}^n n_j^*\right) \cdot \left(1 - \sum_{j=1}^m \lambda_j (l_j^{\prime}+1)\left(l_j + \frac{l_j^{\prime}}{2} \right)\right) 
%         = \sum_{j=1}^m \lambda_j(l_j^{\prime}+1)
% \end{equation}
The equilibrium throughput benchmarks all non-predictive policies within the policy class \(\Pi\). The following proposition demonstrates that the expected throughput of any policy in \(\Pi\) cannot exceed the equilibrium throughput \(\throughput^*\), as specified in Equation \eqref{throughput_fluid_multiple}. 
% Detailed proofs are presented in Appendix~\ref{appendix::proof online policy expected throughput upper bound, multiple type}.

\begin{proposition}\label{prop::online policy expected throughput upper bound, multiple type}
Suppose that the memory capacity satisfies \(C \geq M^*\), where \(M^*\) denotes the equilibrium memory given in \eqref{eq:memory_multi2}. Then, for any online policy \(\pi \in \Pi\), the expected throughput satisfies
\[
\mathbb{E}[\throughput^{(T,\pi)}] \leq \throughput^*,
\]
where \(\throughput^*\) is given in \eqref{throughput_fluid_multiple}.
\end{proposition}
This result has two implications for algorithm design. First, it quantifies the theoretical limit: no non-predictive policy can exceed $\throughput^* = \sum_{j=1}^m \lambda_j l_j'$ when $C \geq M^*$. This establishes a performance benchmark for any scheduling policy. Second, it suggests a design principle: any policy that prevents evictions and maintains balanced stages can approach this throughput. This motivates our WAIT algorithm (Section~\ref{sec:wait}), which uses thresholds to keep the system near equilibrium, preventing the cascading failures that cause FCFS to fall short (see Example~\ref{ex:fcfs_cascade}).

The equilibrium analysis reveals a critical design principle: maintain balanced prompt distribution across stages. Having exactly $n_j^*/(l_j'+1)$ prompts at each stage prevents two failure modes. First, if prompts accumulate at later decode stages---where each consumes more memory---total usage may exceed capacity $C$, forcing evictions that waste partial computation. Second, if prompts concentrate at early stages while later stages starve, the system leaves processing capacity unused, reducing throughput.

This load-balancing principle directly motivates our WAIT algorithm (Section~\ref{sec:wait}). By controlling when prompts advance to the next stage via thresholds, WAIT keeps the system close to the balanced equilibrium, preventing the cascading evictions that plague FCFS (see Example~\ref{ex:fcfs_cascade} in Section~\ref{sec:wait}). The fluid model thus provides both a performance benchmark and a design guide.

In classical queueing theory, fluid models typically lead to ordinary differential equations (ODEs) describing the continuous transition of queue lengths over time. Constructing such ODEs for LLM inference, however, presents technical challenges. Eviction events---triggered when memory usage exceeds capacity---cause discontinuous jumps in the system state, as KV caches of evicted prompts are cleared entirely rather than decreasing smoothly. At boundaries where certain stage populations reach zero or where drift terms vanish, the direction of state transition becomes ambiguous under standard ODE analysis. We therefore adopt a discrete, iteration-based framework that directly characterizes equilibrium and uses it as a performance benchmark pending a complete characterization of such an asymptotic regime.




\begin{comment}Similar to the discussion in the single-type fluid model, when the available memory is larger than the equilibrium memory \( M^* \), we can form a batch with a larger batch size, resulting in \(\text{Arrival} < \text{Completion}\). For example, consider the type \( j \) with the smallest value of \((l_j^\prime+1)\left(l_j+\frac{l_j^\prime}{2}\right)\). Without loss of generality, we set \( j=1 \) and the number of prompts of each stage of type $j=1$ is that $\frac{n_1^*}{l_1^\prime+1}+1$. And the memory \( M^+ \) is defined as
$$M^+ = M^* + (l_1^\prime+1)(l_1+\frac{l_1^\prime}{2})$$
So the arrival and completion ratio is that:
\begin{align*}
\frac{ \text{Arrival}}{ \text{Completion} }=& 
\frac{\left(d_0+d_1M^*+d_1(l_1^\prime+1)\left( l_1 + \frac{l_1^\prime}{2} \right)\right)\sum_{j=1}^m \lambda_j}{ \sum_{j=1}^m  \frac{n_j^*}{l_j^\prime+1} +1}\\
=&\frac{ \sum_{j=1}^m  \frac{n_j^*}{l_j^\prime+1} +d_1(l_1^\prime+1)\left( l_1 + \frac{l_1^\prime}{2} \right)\sum_{j=1}^m\lambda_j }{ \sum_{j=1}^m  \frac{n_j^*}{l_j^\prime+1} +1}<1
\end{align*}
The last inequality comes from that: 
$$ d_1(l_1^\prime+1)\left(l_1+\frac{l_1^\prime}{2}\right)\sum_{j=1}^m \lambda_j  \leq d_1\sum_{j=1}^m \lambda_j (l_j^\prime+1)\left(l_j+\frac{l_j^\prime}{2}\right)<1$$
Similarly, when the the available memory is smaller than the equilibrium memory $M^*$, we can form a batch with a smaller batch size, resulting in \(\text{Arrival} > \text{Completion}\). For example, consider the number of prompts of each stage of type \( j=1 \) is that $\frac{n_1^*}{l_1^\prime+1}-1$. And the memory $M^-$ is given by
$$M^- = M^* - (l_1^\prime+1)(l_1+\frac{l_1^\prime}{2})$$
So the arrival and completion ratio is that:
\begin{align*}
\frac{ \text{Arrival}}{ \text{Completion} }=& 
\frac{\left(d_0+d_1M^*-d_1(l_1^\prime+1)\left( l_1 + \frac{l_1^\prime}{2} \right)\right)\sum_{j=1}^m \lambda_j}{ \sum_{j=1}^m  \frac{n_j^*}{l_j^\prime+1} -1}\\
=&\frac{ \sum_{j=1}^m  \frac{n_j^*}{l_j^\prime+1} -d_1(l_1^\prime+1)\left( l_1 + \frac{l_1^\prime}{2} \right)\sum_{j=1}^m\lambda_j }{ \sum_{j=1}^m  \frac{n_j^*}{l_j^\prime+1} -1}>1
\end{align*}
\end{comment}

% with:
% \[
% n_j^* = \lambda_j (1 + l_j') \quad \forall j.
% \]
% Thus:
% \[
% M^* = \sum_{j=1}^m \lambda_j (1 + l_j') \left( l_j + \frac{l_j' + 1}{2} \right)= C,
% \]
% where the last equality holds since the system dynamics are in equilibrium. Aggregating the results above, we get:
% \[
% \throughput^* = \frac{\sum_{j=1}^mn_j^*}{\sum_{j=1}^m\brk{d_0+d_1n_j^*\prn{l_j+\frac{l_j'}{2}}}}.
% \]

% Now we step We refine the fluid model to incorporate this, defining:
% \begin{itemize}
%     \item \(n_{j, \text{prefill}}^*\): steady-state number of type-\(j\) prompts in the prefill stage per batch.
%     \item \(n_{j, \text{decode}}^*\): steady-state number of type-\(j\) prompts in decode stages per batch.
% \end{itemize}

% The \textit{iteration time} \(\tau^*\), the duration to process one batch, is:
% \[
% \tau^* = \max \left\{ d_f \cdot \max_j \{ l_j \}, \; d_1 \cdot \sum_{j=1}^m n_{j, \text{decode}}^* \left( l_j + \frac{l_j' + 1}{2} \right) \right\},
% \]
% where \(d_f\) and \(d_1\) are prefill and decode processing times per token, and \(\frac{l_j' + 1}{2}\) approximates the average decode-stage memory.

% \textbf{Equilibrium Conditions}: The batch processing rate matches arrivals:
% \[
% \frac{n_{j, \text{prefill}}^*}{\tau^*} = \lambda_j \quad \text{and} \quad \frac{n_{j, \text{decode}}^*}{\tau^*} = \lambda_j l_j' \quad \forall j,
% \]
% reflecting one prefill and \(l_j'\) decode iterations per prompt. Memory must satisfy:
% \[
% \sum_{j=1}^m \left( n_{j, \text{prefill}}^* l_j + n_{j, \text{decode}}^* \left( l_j + \frac{l_j' + 1}{2} \right) \right) \leq C,
% \]
% where \(C\) is the GPU memory capacity.

% \textbf{Throughput}: The equilibrium throughput remains:
% \[
% \throughput^* = \sum_{j=1}^m \lambda_j l_j',
% \]
% consistent with token generation rates.

% This refinement aligns the model with practical constraints, assuming a steady batch composition in equilibrium.

% \subsection{Equilibrium Throughput as a Benchmark}


% \arc{Add details on : 1. overloaded system and underloaded system 2. Lemma on upper bound of fluid throughput over all online non-preemptive policy $\Pi$. 3. Check the calculation.}